% =====================================================================
% Rozdział 1 – Wprowadzenie (~3 str.)
% =====================================================================
\chapter{Wprowadzenie}
\label{ch:wprowadzenie}

% ---------------------------------------------------------------------
\section{Motywacja i~kontekst pracy}
\label{sec:motywacja}
% - Potrzeba percepcji dookólnej (360°) w robotyce mobilnej: robot musi
%   wiedzieć, co jest wokół niego, nie tylko przed nim.
% - Typowe podejścia: obrotowy LiDAR, kamera 360°/rybie oko, kilka kamer
%   na jednym mocnym komputerze – zalety i koszty (cena, pobór mocy,
%   obciążenie komputera).
% - Idea edge AI: przetwarzanie obrazu tam, gdzie powstaje; do komputera
%   trafia kilka bajtów wyniku zamiast strumienia wideo.
% - Dlaczego tanie mikrokontrolery z kamerą: niski koszt węzła, łatwa
%   skalowalność liczby kamer, odciążenie komputera i sieci.
% - Krótko: modele typu FOMO umożliwiają detekcję na mikrokontrolerze
%   z rozsądną liczbą klatek na sekundę.
% - Odniesienie do wszystkich elementów tytułu (wielokamerowy, dookólny,
%   detekcja, czas rzeczywisty).
% Stan: ✅ (tekst do napisania, bez pomiarów)


% ---------------------------------------------------------------------
\section{Cel, zakres i~założenia pracy}
\label{sec:cel-zakres}
% - Cel główny jednym zdaniem (parafraza z APD, czas przeszły):
%   zaprojektowanie i wykonanie systemu z 6 węzłów kamerowych, który
%   wykrywa obiekty wokół robota w czasie rzeczywistym i publikuje wyniki
%   w ROS 2.
% - Cele szczegółowe (lista):
%   - wybór platformy sprzętowej węzła i kamer,
%   - projekt obudowy i montaż systemu na robocie,
%   - zebranie zbioru danych i wytrenowanie modelu FOMO,
%   - oprogramowanie węzła (akwizycja, korekcja dystorsji, inferencja,
%     wysyłka UDP),
%   - procedura kalibracji kamer,
%   - algorytm fuzji detekcji z wielu kamer,
%   - integracja z ROS 2,
%   - ocena jakości i wydajności systemu.
% - Zakres: czego nie ma (brak pomiaru odległości, brak planowania ruchu
%   robota, brak śledzenia obiektów – jeśli nie robisz).
% - Założenia:
%   - system wyznacza KIERUNEK (azymut) do obiektów; odległość nie jest
%     mierzona – świadoma decyzja,
%   - 6 węzłów kamerowych dookoła robota,
%   - mała liczba klas (butelka + tło) – zasada działania; inne klasy
%     wymagają tylko nowego modelu,
%   - wymaganie wydajności: kilka–kilkanaście klatek na sekundę na kamerę,
%   - Wi-Fi 2,4 GHz i UDP w jednej sieci lokalnej,
%   - obliczenia ciężkie (fuzja, ROS 2) na komputerze nadrzędnym;
%     na węzłach tylko detekcja.
% - Ograniczenia: wejście modelu 96x96, warunki oświetleniowe, zakres
%   odległości detekcji.
% Stan: ✅


% ---------------------------------------------------------------------
\section{Układ pracy}
\label{sec:uklad-pracy}
% - Po jednym–dwóch zdaniach o każdym rozdziale 2–6
%   (\ref{ch:teoria}, \ref{ch:projekt}, \ref{ch:implementacja},
%   \ref{ch:testy}, \ref{ch:podsumowanie}).
% - Pisać na końcu, gdy reszta jest gotowa.
% Stan: ⏳ (na koniec)

